\documentclass[11pt]{article}
\usepackage{../dad}
\begin{document}

\courseheader
  {Assignment 0: Find Some Useful Data}
  {LMTH 2080 --- Data and Design with Python}
  {Due Monday, September 14, 2026}
  {Eugene Lang College of Liberal Arts}

\begin{keyidea}
Start with something you want to understand, not with a dataset or a
programming tool.

\begin{center}
\large\bfseries
problem $\longrightarrow$ question $\longrightarrow$ evidence
$\longrightarrow$ data
\end{center}
\end{keyidea}

\section{The Assignment}

Identify a problem or question that genuinely interests you and find
\textbf{two plausible sources of data} that could help you investigate it.

Your question can be large or small. The important thing is that you begin
with a question you care about and then ask what evidence would actually help
you investigate it.

\begin{dataprompt}
\textbf{What problem or question are you interested in?}

Write this before you begin searching for data. You may revise the question
after seeing what data is actually available.
\end{dataprompt}

\section{Find Two Data Sources}

Find two different data sources that could plausibly provide evidence related
to your question. A useful source might be a downloadable CSV or spreadsheet,
a government open-data portal, a public API, data published by a research
organization, an archive, or another source that gives you access to the
underlying observations.

An article, chart, dashboard, or report may help you \emph{find} data, but it
is not automatically the data itself. Whenever possible, locate the
underlying dataset.

\section{Investigate Each Source}

For \textbf{each} source, answer:

\begin{enumerate}[leftmargin=*,itemsep=8pt]
  \item Where does the data come from?
  \item Who produced or collected it?
  \item What does one observation or row represent?
  \item What are several important variables?
  \item What could this data help you understand about your question?
  \item What is missing, uncertain, or limited about the data?
\end{enumerate}

\begin{designprompt}
Imagine someone wanted to use this data to make a decision. What would you
want them to understand about where the data came from and what it fails to
capture?
\end{designprompt}

\section{Choose}

Choose the source that seems more useful for your question and explain why.
Your choice does not have to be the larger, cleaner, or more sophisticated
dataset. A small dataset that closely matches your question may be more useful
than a massive dataset that does not.

If possible, download the data you choose and inspect the actual file.

\section{Create Your README}

Create a file named \texttt{README.md} using this structure:

\begin{verbatim}
# Useful Data

## Question

What problem or question are you interested in?

## Data Source 1

**Source:**
**Produced by:**
**What does one row represent?**
**Important variables:**

### What could this data help you understand?

### What is missing or limited about this data?

## Data Source 2

**Source:**
**Produced by:**
**What does one row represent?**
**Important variables:**

### What could this data help you understand?

### What is missing or limited about this data?

## My Choice

Which dataset seems more useful for your question, and why?
\end{verbatim}

You may write informally, but your README should make it possible for someone
else to understand what you found and why you think it could be useful.

\section{Include the Data When Reasonable}

If your chosen dataset is a reasonably sized downloadable file, include it in
your project directory.

\begin{verbatim}
assignment00/
|-- README.md
`-- data/
    `-- my_data.csv
\end{verbatim}

\begin{checkpoint}
\textbf{Do not put a huge dataset in your Git repository.}

If the data is very large, generated dynamically, accessed through an API, or
cannot reasonably be included, provide a clear link to the source in your
README instead. If useful, you may include a small sample.
\end{checkpoint}

\section{Record Your Work with Git}

Begin by asking what has changed:

\begin{terminal}
git status
\end{terminal}

Stage your README:

\begin{terminal}
git add README.md
\end{terminal}

If you are including a small data file, stage that file as well. Then record
the version:

\begin{terminal}
git commit -m "Complete useful data investigation"
\end{terminal}

Finally:

\begin{terminal}
git status
\end{terminal}

Your instructor will provide any additional GitHub submission instructions.

\section{Generative AI}

For this assignment, \textbf{do not use generative AI to choose your
question, locate your datasets, or evaluate whether the data is useful}.

The searching, inspecting, comparing, and deciding are the work of this
assignment. You may use ordinary web search, documentation, data portals,
course materials, and help from the instructor.

\section{What to Submit}

By the beginning of class on \textbf{Monday, September 14}, you should have:

\begin{itemize}[leftmargin=*]
  \item a question or problem you want to investigate;
  \item two plausible data sources;
  \item a completed \texttt{README.md};
  \item a clear explanation of which source you chose and why;
  \item the chosen data or a link to it, as appropriate; and
  \item at least one Git commit recording your work.
\end{itemize}

\begin{reflection}
\begin{center}
\large
\textbf{Is this data useful for what I am actually trying to understand?}
\end{center}
\end{reflection}

\end{document}
